\begin{salyga}
Lygties $x^2 - 5x + 2 = 0$ šaknys yra $x_1$ ir $x_2$. Nesprendę lygties apskaičiuokite reiškinio reikšmę:
    \[ x_1^2 x_2 + x_1 x_2^2 - 2(x_1^2 + x_2^2) \]
\end{salyga}

\begin{atsakymas}
$-32$.
\end{atsakymas}

\begin{sprendimas}
Pagal Vijeto teoremą $x_1+x_2=5$ ir $x_1x_2=2$. Todėl
\[
x_1^2x_2+x_1x_2^2-2(x_1^2+x_2^2)
=x_1x_2(x_1+x_2)-2\bigl((x_1+x_2)^2-2x_1x_2\bigr)
=10-2\cdot21=-32.
\]
\end{sprendimas}
